\section*{Hoofdstuk \ref{ch:b2:hermitian} --- Hermitische vormen}

\begin{solution}{exo:b2:hermitian:1}
$\langle x, y\rangle = \conj{1}\cdot\iu + \conj{\iu}\cdot 1 = \iu -
\iu = 0$: orthogonaal. $\norm x = \norm y = \sqrt{1 + 1} =
\sqrt2$. Orthonormale basis: $\bigl(\frac{1}{\sqrt2}(1, \iu),\;
\frac{1}{\sqrt2}(\iu, 1)\bigr)$ --- het genormaliseerde paar zelf.
\end{solution}

\begin{solution}{exo:b2:hermitian:2}
Eerste: gelijk aan haar geconjugeerde \omterm{def:b2:linalg:transpose}{getransponeerde} (reële
diagonaal, en $\conj{\iu} = -\iu$ verwisseld): \omterm{def:b2:hermitian:adjoint}{hermitisch} (en dus
normaal); niet \omterm{def:b2:hermitian:adjoint}{unitair} ($A^\dagger A \neq I$: de kolommen hebben
geen \omterm{def:b2:nvs:norm}{norm} $1$).

Tweede: $A^\dagger A = \frac12\begin{pmatrix} 1 & -\iu\\ -\iu &
1\end{pmatrix}\begin{pmatrix} 1 & \iu\\ \iu & 1\end{pmatrix} =
\frac12\begin{pmatrix} 2 & 0\\ 0 & 2\end{pmatrix} = I$: \omterm{def:b2:hermitian:adjoint}{unitair} (en
dus normaal); niet \omterm{def:b2:hermitian:adjoint}{hermitisch}.

Derde: $A^\dagger A = E_{22} \neq E_{11} = AA^\dagger$: niet
normaal (dus noch \omterm{def:b2:hermitian:adjoint}{hermitisch} noch \omterm{def:b2:hermitian:adjoint}{unitair}) --- het standaard
nilpotente tegenvoorbeeld.
\end{solution}

\begin{solution}{exo:b2:hermitian:3}
Eenduidigheid: $A = H + \iu K$ met $H^\dagger = H$, $K^\dagger = K$
dwingt $A^\dagger = H - \iu K$ af, dus $H = \frac{A +
A^\dagger}{2}$ en $K = \frac{A - A^\dagger}{2\iu}$; deze formules
zijn \omterm{def:b2:hermitian:adjoint}{hermitisch} (controle: $\bigl(\frac{A -
A^\dagger}{2\iu}\bigr)^\dagger = \frac{A^\dagger - A}{-2\iu} = K$)
en bouwen $A$ opnieuw op: bestaan.

Normaliteit: $A^\dagger A - AA^\dagger = (H - \iu K)(H + \iu K) -
(H + \iu K)(H - \iu K) = 2\iu(HK - KH)$: dit verdwijnt dan en
slechts dan als $HK = KH$.
\end{solution}

\begin{solution}{exo:b2:hermitian:4}
$\chi_A = (X-2)^2 - 1$: \omterm{def:b2:reduction:eigen}{eigenwaarden} $3$ en $1$. \omterm{def:b2:reduction:eigen}{Eigenvectoren},
door rechtstreekse berekening:
\[
A\begin{pmatrix}1\\ -\iu\end{pmatrix}
= \begin{pmatrix} 2 + \iu(-\iu)\\ -\iu + 2(-\iu)\end{pmatrix}
= \begin{pmatrix} 3\\ -3\iu \end{pmatrix}
= 3\begin{pmatrix}1\\ -\iu\end{pmatrix},
\qquad
A\begin{pmatrix}1\\ \iu\end{pmatrix}
= \begin{pmatrix} 2 + \iu\cdot\iu\\ -\iu + 2\iu\end{pmatrix}
= \begin{pmatrix}1\\ \iu\end{pmatrix}.
\]
Orthonormale eigenbasis: $u_1 = \frac{1}{\sqrt2}(1, -\iu)$
(\omterm{def:b2:reduction:eigen}{eigenwaarde} $3$), $u_2 = \frac{1}{\sqrt2}(1, \iu)$ (\omterm{def:b2:reduction:eigen}{eigenwaarde}
$1$); de orthogonaliteit als in \cref{exo:b2:hermitian:1}. Machten,
via de spectrale projecties $A^k = 3^k u_1u_1^\dagger + 1^k\,
u_2u_2^\dagger$:
\[
A^k = U\begin{pmatrix} 3^k & 0\\ 0 & 1\end{pmatrix}U^\dagger
= \frac{3^k}{2}\begin{pmatrix} 1 & \iu\\ -\iu & 1\end{pmatrix}
+ \frac{1}{2}\begin{pmatrix} 1 & -\iu\\ \iu & 1\end{pmatrix}
= \frac12\begin{pmatrix}
3^k + 1 & (3^k - 1)\iu\\
-(3^k-1)\iu & 3^k + 1
\end{pmatrix}.
\]
(Controleer $k = 1$: dit geeft $A$ terug.)
\end{solution}

\begin{solution}{exo:b2:hermitian:5}
\omterm{def:b2:metric:compact}{Compact}: gesloten (origineel van $I$ onder de \omterm{def:b2:metric:continuity}{continue} afbeelding
$U \mapsto U^\dagger U$) en begrensd (de kolommen zijn
eenheidsvectoren: elementen van modulus $\leq 1$) in
$\mathcal{M}_n(\C) \simeq \R^{2n^2}$.

\omterm{def:b2:reduction:eigen}{Eigenwaarden}: $\operatorname{diag}(\lambda, 1, \dots, 1)$ is
\omterm{def:b2:hermitian:adjoint}{unitair} voor elke $\abs\lambda = 1$. \omterm{def:b2:linalg:det}{Determinant}: $\abs{\det U}^2 =
\det U^\dagger \det U = \det(U^\dagger U) = 1$ (met $\det A^\dagger
= \conj{\det A}$): $\det U$ ligt op de eenheidscirkel.
\end{solution}

\begin{solution}{exo:b2:hermitian:6}
Werk volgens de spectraalstelling in een orthonormale eigenbasis
van $H$: alles herleidt tot scalairen $\lambda \in \R$ (de
\omterm{def:b2:reduction:eigen}{eigenwaarden}). $H + \iu I$ heeft \omterm{def:b2:reduction:eigen}{eigenwaarden} $\lambda + \iu \neq
0$: inverteerbaar. $U$ heeft \omterm{def:b2:reduction:eigen}{eigenwaarden} $\mu = \frac{\lambda -
\iu}{\lambda + \iu}$, van modulus $1$ ($\abs{\lambda - \iu} =
\abs{\lambda + \iu}$ voor reële $\lambda$): $U^\dagger U = I$ geldt
omdat $U$ \omterm{def:b2:hermitian:adjoint}{unitair} \omterm{def:b2:reduction:diag}{diagonaliseerbaar} is met \omterm{def:b2:reduction:eigen}{eigenwaarden} van modulus
$1$ (zij is diagonaal in de gekozen orthonormale basis). En $\mu =
1$ zou $-\iu = \iu$ afdwingen: onmogelijk, dus $1 \notin
\operatorname{Sp} U$. (De Cayley-transformatie beeldt het
\omterm{def:b2:hermitian:adjoint}{hermitische} geval af op het \omterm{def:b2:hermitian:adjoint}{unitaire} min één punt --- de
matrixversie van de afbeelding van $\R$ naar de cirkel.)
\end{solution}

\begin{solution}{exo:b2:hermitian:7}
Voor een eigenpaar $Ax = \lambda x$ ($x \neq 0$) is
$\lambda\norm x^2 = \langle x, Ax\rangle > 0$, dus $\lambda > 0$
(reëel was zij al, \cref{prop:b2:hermitian:eigenvalues}).
Vierkantswortel: in een spectrale basis is $B =
U\operatorname{diag}(\sqrt{\lambda_i})U^\dagger$: \omterm{def:b2:hermitian:adjoint}{hermitisch},
positief definiet, $B^2 = A$. \omterm{def:b2:linalg:det}{Determinant}: het product van de
positieve \omterm{def:b2:reduction:eigen}{eigenwaarden}.
\end{solution}

\begin{solution}{exo:b2:hermitian:8}
\begin{enumerate}
  \item $\norm{u(x)}^2 = \langle u(x), u(x)\rangle = \langle x,
        u^*u(x)\rangle = \langle x, uu^*(x)\rangle =
        \norm{u^*(x)}^2$. Toegepast op de normale $u -
        \lambda\,\mathrm{id}$ (haar \omterm{def:b2:quadratic:adjoint}{adjunct} is $u^* -
        \conj\lambda$, en de normaliteit wordt geërfd):
        $\norm{(u - \lambda)x} = \norm{(u^* - \conj\lambda)x}$,
        dus de kernen vallen samen.
  \item Voor \omterm{def:b2:reduction:eigen}{eigenvectoren} $u(x) = \lambda x$, $u(y) = \mu y$
        ($\lambda \neq \mu$) geeft (1) dat $u^*(x) = \conj\lambda
        x$; dan is
        \[
        \lambda\langle y, x\rangle = \langle y, u(x)\rangle
        = \langle u^*(y), x\rangle = \langle \conj\mu\, y, x\rangle
        = \mu \langle y, x\rangle ,
        \]
        dus $\langle y, x\rangle = 0$. Stabiliteit van
        $E_\lambda^\perp$: voor $x \perp E_\lambda$ en $z \in
        E_\lambda$ is $\langle z, u(x)\rangle = \langle u^*(z),
        x\rangle = \langle\conj\lambda z, x\rangle = 0$.
  \item Inductie naar de dimensie: over $\C$ heeft $u$ een
        \omterm{def:b2:reduction:eigen}{eigenvector} $e_1$ (normaliseer haar); haar orthogonale
        complement is stabiel onder $u$ (volgens (2)) \emph{en}
        onder $u^*$ (hetzelfde argument met verwisselde rollen),
        dus de beperking is normaal: pas inductie toe en plak de
        orthonormale eigenbasissen aaneen. Omgekeerd voldoet een
        \omterm{def:b2:hermitian:adjoint}{unitair} \omterm{def:b2:reduction:diag}{diagonaliseerbare} $u = UDU^\dagger$ aan $u^*u
        = U\conj D D U^\dagger = U D\conj D U^\dagger = uu^*$:
        normaal.
\end{enumerate}
\end{solution}

\begin{solution}{exo:b2:hermitian:9}
\omterm{def:b2:reduction:eigen}{Eigenwaarden}: voor $u(x) = \lambda x$, $x \neq 0$:
\[
\lambda\norm x^2 = \langle x, u(x)\rangle
= \langle u^*(x), x\rangle = -\langle u(x), x\rangle
= -\conj{\langle x, u(x)\rangle} = -\conj\lambda\,\norm x^2 ,
\]
dus $\lambda = -\conj\lambda$: zuiver imaginair. Omdat $(\iu u)^* =
-\iu\,u^*$ (de \omterm{def:b2:quadratic:adjoint}{adjunct} is geconjugeerd lineair in de scalairen),
geeft $u^* = u$ dat $(\iu u)^* = -\iu u$: de afbeelding $u \mapsto
\iu u$ zendt \omterm{def:b2:hermitian:adjoint}{hermitisch} naar \omterm{def:b2:hermitian:adjoint}{antihermitisch}, met inverse $w \mapsto
-\iu w$: een bijectie. Een \omterm{def:b2:hermitian:adjoint}{antihermitische} $u$ voldoet aan $u^*u =
-u^2 = uu^*$: normaal, en dus \omterm{def:b2:hermitian:adjoint}{unitair} \omterm{def:b2:reduction:diag}{diagonaliseerbaar} volgens
\cref{exo:b2:hermitian:8}.
\end{solution}

\begin{solution}{exo:b2:hermitian:10}
\begin{enumerate}
  \item $S$ permuteert een orthonormale basis: $\norm{Sx} = \norm
        x$, dus $S$ is \omterm{def:b2:hermitian:adjoint}{unitair}. Indexeren we de coördinaten met
        $j = 0, \dots, n-1$ modulo $n$, dan is $(Sx)_j =
        x_{j-1}$, dus voor $(f_k)_j = \frac{\omega^{jk}}{\sqrt
        n}$:
        \[
        (Sf_k)_j = \frac{\omega^{(j-1)k}}{\sqrt n}
        = \omega^{-k}\,(f_k)_j :
        \qquad Sf_k = \omega^{-k}f_k .
        \]
        Orthonormaliteit: $\langle f_k, f_l\rangle = \frac1n
        \sum_j \omega^{j(l-k)} = \delta_{kl}$ (de meetkundige som
        van een niet-triviale eenheidswortel verdwijnt).
  \item $Cf_k = \sum_j c_j S^jf_k = \bigl(\sum_j
        c_j\omega^{-jk}\bigr)f_k = \widehat c(k)\,f_k$: elke
        circulante matrix is diagonaal in de orthonormale
        Fourierbasis, en dus normaal, met spectrum $\{\widehat
        c(k)\}$. (De basisovergang is de discrete
        Fouriertransformatie: convolutie wordt
        vermenigvuldiging.)
\end{enumerate}
\end{solution}

\begin{solution}{exo:b2:hermitian:11}
($\Leftarrow$) Zij $P^2 = P = P^*$. Elke $v$ splitst als $v = Pv +
(v - Pv)$ met $Pv \in \operatorname{im} P$ en $P(v - Pv) = 0$. De
twee stukken zijn orthogonaal: voor alle $x, y$ is
\[
\langle Px, (I - P)y\rangle = \langle x, P(I-P)y\rangle
= \langle x, (P - P^2)y\rangle = 0 :
\]
$\ker P \perp \operatorname{im} P$, dus $P$ is de orthogonale
projectie op haar beeld. ($\Rightarrow$) Is $P$ de orthogonale
projectie op $F = \operatorname{im}P$, dan is voor alle $x, y$
$\langle Px, y\rangle = \langle Px, Py\rangle$ (de component $y -
Py \perp F$ valt weg) en \omterm{def:b2:quadratic:adjoint}{symmetrisch} $\langle x, Py\rangle =
\langle Px, Py\rangle$: dus $\langle Px, y\rangle = \langle x,
Py\rangle$, dat wil zeggen $P^* = P$. Matrices: op $\C v$ ($\norm v
= 1$): $Px = v\,\langle v, x\rangle$, dat wil zeggen $P =
vv^\dagger$; op $\operatorname{Vect}(v_1, \dots, v_k)$
orthonormaal: $P = \sum_i v_iv_i^\dagger$.
\end{solution}

\begin{solution}{exo:b2:hermitian:12}
Diagonaliseer $A = U D U^\dagger$ (spectraalstelling), met $D$
diagonaal met elementen uit de $\lambda_i$. Dan is $P_i = L_i(A) =
U L_i(D)U^\dagger$, en $L_i(D)$ is diagonaal met elementen
$L_i(\lambda_j) = \delta_{ij}$: enen precies op de plaatsen van
$E_i$. Dus $P_i$ is \omterm{def:b2:hermitian:adjoint}{hermitisch} ($L_i$ reëel, $D$ reëel),
idempotent, met beeld $E_i$ en kern $\bigoplus_{j\neq i}E_j =
E_i^\perp$ (orthogonaliteit van de \omterm{def:b2:reduction:eigen}{eigenruimten}): de orthogonale
projectie op $E_i$ (\cref{exo:b2:hermitian:11}). Disjuncte
diagonaalpatronen geven $P_iP_j = 0$ ($i \neq j$); $\sum_i L_i = 1$
(graad $< p$, waarde $1$ in $p$ punten), dus $\sum P_i = I$; en
$\sum_i \lambda_iL_i(\lambda_j) = \lambda_j$ geeft $A = \sum
\lambda_iP_i$. Voor een willekeurige veelterm $f$ heeft $f(D)$ de
diagonaal $f(\lambda_j)$, dus
\[
f(A) = \sum_{i=1}^{p} f(\lambda_i)\,P_i :
\]
functies van $A$ worden spectraal berekend --- de rekenkunde die
het volume van bachelorjaar 3 uitbreidt tot \omterm{def:b2:metric:continuity}{continue} $f$ en verder.
\end{solution}

\begin{solution}{pb:b2:hermitian:1}
\textbf{1.} $\conj{\langle x, Ax\rangle} = \langle Ax, x\rangle
= \langle x, A^*x\rangle = \langle x, Ax\rangle$: reëel. Schrijven
we $x = \sum c_ie_i$, dan is
\[
R_A(x) = \frac{\sum_i\lambda_i\abs{c_i}^2}
{\sum_i\abs{c_i}^2} ,
\]
een gewogen gemiddelde van de \omterm{def:b2:reduction:eigen}{eigenwaarden}: het ligt in
$\intcc{\lambda_n}{\lambda_1}$, met de grenzen aangenomen in $e_1$
en $e_n$.

\textbf{2.} Werk voor reële $t$ en willekeurige $v$ uit dat $R_A(x
+ tv) = \frac{N(t)}{D(t)}$ met
\[
N(t) = \langle x, Ax\rangle + 2t\Re\langle v, Ax\rangle +
t^2\langle v, Av\rangle,
\quad
D(t) = \norm x^2 + 2t\Re\langle v, x\rangle + t^2\norm v^2 .
\]
De afgeleide in $t = 0$ is
\[
\frac{2}{\norm x^2}\,
\Re\bigl\langle v,\ Ax - R_A(x)\,x\bigr\rangle .
\]
Zij verdwijnt voor alle $v$ dan en slechts dan als $\Re\langle v,
w\rangle = 0$ voor alle $v$, met $w = Ax - R_A(x)x$; $v$ door $\iu
v$ vervangen doodt ook het imaginaire deel: $w = 0$, dat wil zeggen
$Ax = R_A(x)x$. De kritieke punten van $R_A$ zijn dus precies de
\omterm{def:b2:reduction:eigen}{eigenvectoren}, met als kritieke waarde de \omterm{def:b2:reduction:eigen}{eigenwaarde}.

\textbf{3.} Voor $x = \sum_{i\leq k}c_ie_i \in V_k$ is $R_A(x)$
een gewogen gemiddelde van $\lambda_1, \dots, \lambda_k$, en dus
$\geq \lambda_k$, met gelijkheid in $e_k$: $\min_{V_k} R_A =
\lambda_k$. \omterm{def:b2:quadratic:adjoint}{Symmetrisch} betreft het gemiddelde op $W_k$ de waarden
$\lambda_k, \dots, \lambda_n$: $\max_{W_k}R_A = \lambda_k$.

\textbf{4.} Zij $\dim V = k$. Dan is $\dim V + \dim W_k = n + 1 >
n$, dus er is een eenheidsvector $x \in V \cap W_k$, en $R_A(x)
\leq \lambda_k$ (vraag 3): $\min_{V}R_A \leq \lambda_k$ voor elke
zulke $V$. Omdat $V_k$ de waarde $\lambda_k$ aanneemt, is het
max-min gelijk aan $\lambda_k$. De min-maxformule is hetzelfde
argument met omgekeerde rollen ($\dim W = n - k + 1$ dwingt $W \cap
V_k \neq \{0\}$ af, dus $\max_W R_A \geq \lambda_k$, aangenomen in
$W_k$).

\textbf{5.} $R_B(x) = R_A(x) + \frac{\langle x, (B-A)x\rangle}
{\norm x^2} \geq R_A(x)$ puntsgewijs. Nemen we het minimum over
een willekeurige $k$-dimensionale $V$ en daarna het maximum over
$V$, dan geeft vraag 4 dat $\lambda_k(B) \geq \lambda_k(A)$.

\textbf{6.} Grassmann: $\dim(V\cap W) = \dim V + \dim W - \dim(V +
W) \geq \dim V + \dim W - n > 0$. Tweemaal toegepast:
\[
\dim(V_1\cap V_2\cap V_3)
\geq \dim(V_1\cap V_2) + \dim V_3 - n
\geq \dim V_1 + \dim V_2 + \dim V_3 - 2n .
\]

\textbf{7.} De deelruimten $W_i(A)$, $W_j(B)$ (vraag 3, voor $A$ en
$B$) en $V_{i+j-1}(A+B)$ hebben dimensies $(n-i+1) + (n-j+1) +
(i+j-1) = 2n + 1 > 2n$: volgens vraag 6 is er een eenheidsvector
$x$ in alle drie. Dan is
\[
\lambda_{i+j-1}(A+B) \leq R_{A+B}(x)
= R_A(x) + R_B(x) \leq \lambda_i(A) + \lambda_j(B),
\]
de linkerongelijkheid omdat $x \in V_{i+j-1}(A+B)$ (vraag 3), de
rechter wegens de twee $W$'s.

\textbf{8.} In een spectrale basis van $E$: $\norm{Ex}^2 = \sum
\lambda_k(E)^2\abs{c_k}^2 \leq \max_k\lambda_k(E)^2\,\norm x^2$,
aangenomen in de bijbehorende \omterm{def:b2:reduction:eigen}{eigenvector}: $\vertiii E_2 =
\max_k\abs{\lambda_k(E)}$. Weyl met $j = 1$: $\lambda_k(A+E) \leq
\lambda_k(A) + \lambda_1(E) \leq \lambda_k(A) + \vertiii E_2$; dit
toegepast op $(A+E) + (-E)$ geeft $\lambda_k(A) \leq \lambda_k(A+E)
+ \vertiii E_2$. Samen: $\abs{\lambda_k(A+E) - \lambda_k(A)} \leq
\vertiii E_2$.

\textbf{9.} Ondergrenzen: $P \geq 0$ en vraag 5. Bovengrens: $P$
heeft rang $1$, dus $\lambda_2(P) = 0$; Weyl met $i = k-1$, $j =
2$:
\[
\lambda_k(A + P) \leq \lambda_{k-1}(A) + \lambda_2(P)
= \lambda_{k-1}(A) .
\]

\textbf{10.} $A + E = \begin{pmatrix} 2 & 1+\iu\\ 1-\iu &
2\end{pmatrix}$: \omterm{def:b2:reduction:charpoly}{karakteristieke veelterm} $(2-\lambda)^2 -
\abs{1+\iu}^2 = (2-\lambda)^2 - 2$, spectrum $\{2 + \sqrt2,\ 2 -
\sqrt2\}$. Tegenover $\operatorname{Sp}A = \{3, 1\}$:
\[
\abs{(2+\sqrt2) - 3} = \abs{(2-\sqrt2) - 1} = \sqrt2 - 1
\approx 0.414 \leq 1 = \vertiii E_2 . \checkmark
\]

\textbf{11.} Beschouw $\C^{n-1} = \operatorname{Vect}(e_1, \dots,
e_{n-1})$ binnen $\C^n$ (standaardbasis): voor $x$ daarin is
$\langle x, Bx\rangle = \langle x, Ax\rangle$, dus $R_B$ is de
beperking van $R_A$. Bovengrens: het max-min voor $\lambda_k(B)$
loopt over de $k$-dimensionale deelruimten \emph{van} $\C^{n-1}$,
een deelfamilie van die van $\C^n$: $\lambda_k(B) \leq
\lambda_k(A)$. Ondergrens: het min-max voor $\lambda_k(B)$ loopt
over deelruimten van $\C^{n-1}$ van dimensie $(n-1)-k+1 = n-k$;
elk daarvan is ook een deelruimte van $\C^n$ van dimensie $n -
(k+1) + 1$, dus haar maximum is $\geq \lambda_{k+1}(A)$:
$\lambda_k(B) \geq \lambda_{k+1}(A)$.

\textbf{12.} Verwijder de rijen en kolommen één voor één en rijg
vraag 11 aaneen: elke verwijdering verschuift de onderste index met
één, wat $\lambda_{k+m}(A) \leq \lambda_k(B) \leq \lambda_k(A)$
geeft.

\textbf{13.} Conjugeren van $A$ met een permutatiematrix (\omterm{def:b2:hermitian:adjoint}{unitair})
verandert noch het spectrum noch de multiverzameling van
diagonaalelementen: neem dus aan dat $d_1, \dots, d_k$ de leidende
posities innemen. De leidende hoofdondermatrix $B$ van formaat
$k\times k$ heeft dan $\operatorname{tr} B = \sum_{i\leq k}d_i$, en
haar \omterm{def:b2:reduction:eigen}{eigenwaarden} voldoen aan $\mu_i(B) \leq \lambda_i(A)$ (vraag
12): sommeren geeft $\sum_{i\leq k}d_i \leq \sum_{i\leq
k}\lambda_i(A)$. Bij $k = n$ zijn beide leden $\operatorname{tr}
A$.

\textbf{14.} De keuze $x_i = e_i$ geeft de waarde $\sum_{i\leq
k}\lambda_i$: het maximum is $\geq$. Omgekeerd laat een
orthonormale familie $(x_1, \dots, x_k)$ zich uitbreiden tot een
orthonormale basis, dat wil zeggen tot een \omterm{def:b2:hermitian:adjoint}{unitaire} $U$ met de
$x_i$ als eerste kolommen; dan is $\sum_i\langle x_i, Ax_i\rangle$
de som van de eerste $k$ diagonaalelementen van $U^\dagger AU$, die
volgens vraag 13 hoogstens de som van haar $k$ grootste
\omterm{def:b2:reduction:eigen}{eigenwaarden} is --- namelijk $\sum_{i\leq k}\lambda_i(A)$. Het
maximumprincipe van Ky Fan volgt.

\textbf{15.} Interlacering ($\operatorname{Sp}A = \{2+\sqrt2, 2,
2-\sqrt2\}$, $\operatorname{Sp}B = \{3, 1\}$):
\[
2 \leq 3 \leq 2 + \sqrt2,
\qquad
2 - \sqrt2 \leq 1 \leq 2 . \checkmark
\]
Schur met diagonaal $(2,2,2)$: $2 \leq 2+\sqrt2$; $4 \leq 4 +
\sqrt2$; $6 = 6$ (het spoor). \checkmark

\textbf{16.} $b_{ii} - a_{ii} = \langle e_i, (B-A)e_i\rangle \geq
0$; sommeren geeft de sporen. Voor elke $C$: $\langle x,
C^\dagger(B - A)Cx\rangle = \langle Cx, (B-A)(Cx)\rangle \geq 0$:
$C^\dagger AC \leq C^\dagger BC$.

\textbf{17.} $A \geq 0$ is duidelijk; $B - A = \begin{pmatrix} 1 &
1\\ 1 & 1\end{pmatrix}$ is positief semidefiniet (\omterm{def:b2:reduction:eigen}{eigenwaarden} $2,
0$): $A \leq B$. Maar
\[
B^2 = \begin{pmatrix} 5 & 3\\ 3 & 2\end{pmatrix},
\qquad
B^2 - A^2 = \begin{pmatrix} 4 & 3\\ 3 & 2\end{pmatrix},
\qquad \det(B^2 - A^2) = -1 < 0 :
\]
niet positief semidefiniet. Kwadrateren eerbiedigt de ordening van
Loewner niet.

\textbf{18.} Zij $S = \sqrt A$ en $T = \sqrt B$ (\omterm{def:b2:hermitian:adjoint}{hermitisch}
positief semidefiniet; \cref{exo:b2:hermitian:7} met dezelfde
spectrale formule uitgebreid tot het semidefiniete geval). $T - S$
is \omterm{def:b2:hermitian:adjoint}{hermitisch}; zij $\mu$ een willekeurige \omterm{def:b2:reduction:eigen}{eigenwaarde} en $v$ een
\omterm{def:b2:reduction:eigen}{eigenvector} van \omterm{def:b2:nvs:norm}{norm} $1$. Uit $T^2 - S^2 = T(T - S) + (T - S)S$
volgt
\[
0 \leq \langle v, (B - A)v\rangle
= \langle Tv, (T-S)v\rangle + \langle (T-S)v, Sv\rangle
= \mu\bigl(\langle v, Tv\rangle + \langle v,
Sv\rangle\bigr)
\]
($\mu$ is reëel). Is $\langle v, Tv\rangle + \langle v, Sv\rangle >
0$, dan is $\mu \geq 0$. Verdwijnt zij, dan verdwijnen beide
niet-negatieve termen; $\langle v, Tv\rangle = \norm{T^{1/2}v}^2 =
0$ dwingt $Tv = 0$ af, en evenzo $Sv = 0$, dus $\mu v = (T - S)v =
0$ en $\mu = 0$. Alle \omterm{def:b2:reduction:eigen}{eigenwaarden} van $T - S$ zijn $\geq 0$:
$\sqrt A \leq \sqrt B$.

\textbf{19.} Congruentie met $A^{-1/2}$ (vraag 16): $I \leq M :=
A^{-1/2}BA^{-1/2}$. Dus zijn alle \omterm{def:b2:reduction:eigen}{eigenwaarden} van $M$ groter dan
of gelijk aan $1$, en die van $M^{-1}$ liggen in $\intoc{0}{1}$:
$M^{-1} \leq I$. Maar $M^{-1} = A^{1/2}B^{-1}A^{1/2}$; $M^{-1} \leq
I$ congrueren met $A^{-1/2}$ geeft $B^{-1} \leq A^{-1}$.

\textbf{20.} $\sqrt A^{-1}(AB)\sqrt A = \sqrt A\,B\sqrt A$: dus is
$AB$ gelijkvormig met de \omterm{def:b2:hermitian:adjoint}{hermitische} positief definiete $\sqrt
A\,B\sqrt A$ (definiet: $\langle x, \sqrt AB\sqrt Ax\rangle =
\langle \sqrt Ax, B\sqrt Ax\rangle > 0$): haar \omterm{def:b2:reduction:eigen}{eigenwaarden} zijn
reëel en positief. Bovendien is
\[
\langle x, \sqrt AB\sqrt Ax\rangle
\leq \lambda_{\max}(B)\,\norm{\sqrt Ax}^2
= \lambda_{\max}(B)\,\langle x, Ax\rangle
\leq \lambda_{\max}(A)\lambda_{\max}(B)\norm x^2 ,
\]
dus $\lambda_{\max}(AB) = \max R_{\sqrt AB\sqrt A} \leq
\lambda_{\max}(A)\lambda_{\max}(B)$.

\textbf{21.} Met $v_k = (\sin j\theta_k)_j$ en de identiteit
$\sin((j-1)\theta) + \sin((j+1)\theta) = 2\sin(j\theta)\cos\theta$
geldt voor $2 \leq j \leq n-1$:
\[
(T_nv_k)_j = -\sin((j{-}1)\theta_k) + 2\sin(j\theta_k) -
\sin((j{+}1)\theta_k)
= (2 - 2\cos\theta_k)\sin(j\theta_k) .
\]
Rij $1$ werkt omdat $\sin(0\cdot\theta_k) = 0$, rij $n$ omdat
$\sin((n+1)\theta_k) = \sin(k\pi) = 0$: de randvoorwaarden
selecteren precies $\theta_k = \frac{k\pi}{n+1}$. Dus $T_nv_k =
4\sin^2\bigl(\frac{k\pi}{2(n+1)}\bigr)v_k$; de $n$ waarden zijn
verschillend en liggen in $\intoo04$, en de $v_k \neq 0$: dit is
het hele spectrum, positief, dus $T_n$ is positief definiet.

\textbf{22.} $\min_id_i\,I \leq D \leq \max_id_i\,I$, dus $T_n +
\min_id_i\,I \leq T_n + D \leq T_n + \max_id_i\,I$ (het optellen
van $T_n$ bewaart de ordening); vraag 5 en $\lambda_k(T_n + cI) =
\lambda_k(T_n) + c$ geven de insluiting.

\textbf{23.} $T_2 = \begin{pmatrix} 2 & -1\\ -1 &
2\end{pmatrix}$ heeft spectrum $\{3, 1\}$, en
\[
2 - \sqrt2 \;\leq\; 1 \;\leq\; 2 \;\leq\; 3 \;\leq\; 2 +
\sqrt2 :
\]
de interlacering van Cauchy is nagegaan. (Dit zijn dezelfde spectra
als in vraag 15: conjugeren met $\operatorname{diag}(1,-1,1)$
verandert het teken buiten de diagonaal.) Lezing via grafen: $T_2$
is de matrix van laplaciaantype van het pad waarvan het laatste
punt is geschrapt --- een hoofdondermatrix, precies de situatie van
vraag 11.

\textbf{24.} De extreme \omterm{def:b2:reduction:eigen}{eigenwaarden} gedragen zich als
\[
\lambda_{\min}(T_n) = 4\sin^2\frac{\pi}{2(n+1)}
\sim \frac{\pi^2}{(n+1)^2},
\qquad
\lambda_{\max}(T_n) = 4\cos^2\frac{\pi}{2(n+1)}
\longrightarrow 4 .
\]
Bijgevolg is
\[
\kappa_n = \frac{\lambda_{\max}}{\lambda_{\min}}
\sim \frac{4(n+1)^2}{\pi^2} :
\]
hoe fijner het rooster, hoe slechter geconditioneerd de discrete
tweede afgeleide --- een feit dat het ontwerp van de numerieke
lineaire algebra stuurt.

\textbf{25.} (i) De nulpunten van de \omterm{def:b2:reduction:charpoly}{karakteristieke veelterm}
kunnen bij verstoringen van een algemene matrix wild bewegen, maar
de min-maxkarakterisering klemt elke \omterm{def:b2:hermitian:adjoint}{hermitische} \omterm{def:b2:reduction:eigen}{eigenwaarde} tussen
expliciete optimalisatiewaarden, wat de $1$-lipschitzstabiliteit
van de vragen 8--9 afdwingt. (ii) De vragen 1, 5 en 16--20
gebruikten alleen de extreme Rayleigh-waarden; Weyl, interlacering,
Schur en Ky Fan (vragen 7--14) hadden werkelijk de volledige
min-max over deelruimten nodig. (iii) De ordening van Loewner maakt
van deze scalaire ongelijkheden een rekenkunde van
matrixongelijkheden --- met echte valkuilen (vraag 17) en echte
stellingen (vragen 18--19). (iv) Deze gereedschappen zijn het
dagelijks brood van de numerieke wiskunde (de stijfheidsmatrices
van vraag 24), van de kwantumverstoringstheorie (Weyl: de
energieniveaus verplaatsen zich hoogstens over de \omterm{def:b2:nvs:norm}{norm} van de
verstoring) en van het min-maxprincipe voor \omterm{def:b2:metric:compact}{compacte}
\omterm{def:b2:quadratic:adjoint}{zelftoegevoegde} operatoren in het volume van bachelorjaar 3.
\end{solution}
